\chapter{Il modulo \code{async-queue}}
\label{cap:09-async-queue}

\code{async-queue} \`e il modulo che rende possibile l'invocazione asincrona, e
--- meno ovviamente --- \`e anche l'unico che fornisce alla piattaforma
un'osservazione reale dello stato delle code. Il primo ruolo \`e quello che il
suo nome dichiara; il secondo \`e ci\`o che lo rende una dipendenza per
\code{concurrency-control} e una dipendenza parziale per \code{autoscaler}.

Sono 682 righe di codice in nove classi.

\section{Struttura}

\begin{center}
\small
\begin{adjustbox}{max width=\textwidth}
\begin{tabular}{@{}lp{9cm}@{}}
\toprule
\textbf{Classe} & \textbf{Ruolo} \\
\midrule
\code{FunctionQueueState} & Stato di una singola funzione: coda limitata,
contatore delle richieste in volo, limiti di concorrenza. \\
\code{QueueManager} & Mappa da nome di funzione a stato, registrazione delle
metriche, notifica del lavoro disponibile. \\
\code{Scheduler} & Ciclo a thread singolo che consuma dalle code e dispaccia. \\
\code{QueueBackedEnqueuer} & Implementazione di \code{InvocationEnqueuer}. \\
\code{QueueBackedMetricsSource} & Implementazione di \code{ScalingMetricsSource}. \\
\code{QueueConcurrencyControlMetrics} & Metriche descrittive del controllore
attivo per funzione. \\
\code{WorkSignaler} & Interfaccia a un metodo per la notifica. \\
\bottomrule
\end{tabular}
\end{adjustbox}
\end{center}

\section{Lo stato di una funzione}

\code{FunctionQueueState} incapsula quattro grandezze: la coda dei compiti
(\code{ArrayBlockingQueue} di dimensione fissa), il contatore delle richieste in
volo, il limite di concorrenza \emph{configurato} e il limite di concorrenza
\emph{effettivo}.

La distinzione fra i due limiti \`e la chiave dell'intero modulo. Il limite
configurato \`e ci\`o che l'utente ha dichiarato nella specifica della funzione.
Il limite effettivo \`e ci\`o che il modulo \code{concurrency-control}, se
presente, ha deciso di applicare in questo momento. In assenza di quel modulo i
due coincidono sempre.

\subsection{Acquisizione dello slot}

L'ammissione al dispatch \`e un'operazione atomica di test-e-incremento:

\begin{lstlisting}[language=Java, caption={\code{FunctionQueueState.tryAcquireSlot()}.}]
/**
 * Atomically checks if dispatch is allowed and increments inFlight if so.
 * This prevents race conditions where multiple threads could both pass
 * canDispatch() check and then both increment, exceeding the concurrency limit.
 */
public boolean tryAcquireSlot() {
    while (true) {
        int current = inFlight.get();
        if (current >= effectiveConcurrency) {
            return false;
        }
        if (inFlight.compareAndSet(current, current + 1)) {
            return true;
        }
        // CAS failed, another thread modified - retry
    }
}
\end{lstlisting}

L'alternativa ingenua --- verificare \code{canDispatch()} e poi incrementare ---
\`e corretta solo con un unico consumatore. Il modulo ne ha uno solo oggi (lo
scheduler \`e a thread singolo), ma il metodo \`e esposto anche al percorso
sincrono attraverso \code{InvocationEnqueuer.tryAcquireSlot}, che gira su thread
arbitrari. La primitiva atomica \`e quindi necessaria, non difensiva.

\subsection{Il rilascio con pavimento a zero}

Il rilascio \`e simmetrico ma con una protezione aggiuntiva:

\begin{lstlisting}[language=Java]
private void decrementInFlightNonNegative() {
    while (true) {
        int current = inFlight.get();
        if (current == 0) { return; }
        if (inFlight.compareAndSet(current, current - 1)) { return; }
    }
}
\end{lstlisting}

Un decremento su un contatore gi\`a a zero viene ignorato anzich\'e produrre un
valore negativo. La condizione non dovrebbe verificarsi --- il
\cref{cap:06-nucleo} descrive il meccanismo che rende idempotente il rilascio
per tentativo --- ma se si verificasse, un contatore negativo si propagherebbe
direttamente nella metrica \code{function\_inFlight}, che \`e la sorgente su cui
decidono sia l'autoscaler sia il controllore di concorrenza. Un errore di
contabilit\`a in questo punto non produce un malfunzionamento locale: produce
decisioni di politica sbagliate a valle, per un tempo indeterminato. Il pavimento
a zero \`e la contenzione minima di un errore che si propaga.

\subsection{Interazione fra limite configurato ed effettivo}

I due metodi che scrivono i limiti codificano una semantica che vale la pena
esplicitare:

\begin{lstlisting}[language=Java]
public void concurrency(int concurrency) {
    int previousConfigured = this.configuredConcurrency;
    int normalized = Math.max(1, concurrency);
    this.configuredConcurrency = normalized;
    if (effectiveConcurrency == previousConfigured || effectiveConcurrency > normalized) {
        // Preserve fixed-mode semantics: effective limit tracks configured limit
        // (and clamps down on shrink).
        effectiveConcurrency = normalized;
    }
}

public void setEffectiveConcurrency(int effectiveConcurrency) {
    int clamped = Math.max(1, effectiveConcurrency);
    if (clamped > configuredConcurrency) { clamped = configuredConcurrency; }
    this.effectiveConcurrency = clamped;
}
\end{lstlisting}

Due invarianti emergono. Il primo: \textbf{il limite effettivo non pu\`o mai
superare quello configurato}. Un controllore adattivo pu\`o restringere ma non
allargare oltre ci\`o che l'utente ha autorizzato; la specifica della funzione
resta un tetto. Il secondo: un aggiornamento del limite configurato trascina
quello effettivo \emph{solo} se questo non era stato scostato da un controllore,
oppure se il nuovo tetto \`e pi\`u basso. In pratica: cambiare la configurazione
non cancella una decisione adattiva in corso, ma un abbassamento del tetto la
vincola immediatamente.

\section{Lo scheduler}

Lo scheduler \`e un ciclo su thread singolo. La sua struttura \`e determinata da
un requisito: non deve sondare le code inattive.

\begin{lstlisting}[language=Java]
private void loop() {
    while (running.get()) {
        String functionName = activeFunctions.poll(500, TimeUnit.MILLISECONDS);
        if (functionName != null) {
            enqueuedFunctions.remove(functionName);
            processFunction(functionName);
        }
        ...
    }
}
\end{lstlisting}

La coda \code{activeFunctions} contiene i nomi delle funzioni che hanno lavoro
pendente; l'insieme \code{enqueuedFunctions} evita che lo stesso nome vi compaia
pi\`u volte. Un accodamento notifica lo scheduler solo se il nome non era gi\`a
segnalato. Il risultato \`e che il costo del ciclo \`e proporzionale al numero di
funzioni \emph{attive}, non al numero di funzioni registrate: registrare mille
funzioni inattive non costa nulla.

Il timeout di 500\,ms sulla \code{poll} \`e ci\`o che rende il ciclo terminabile:
senza di esso l'arresto richiederebbe un'interruzione del thread bloccato.

\subsection{Il batch limitato e l'equit\`a}

\begin{lstlisting}[language=Java]
private static final int MAX_BATCH_PER_FUNCTION = 2;

private void processFunction(String functionName) {
    ...
    int dispatched = 0;
    while (running.get() && dispatched < MAX_BATCH_PER_FUNCTION && state.tryAcquireSlot()) {
        InvocationTask task = state.poll();
        if (task == null) { state.releaseSlot(); break; }
        dispatched++;
        SchedulerDispatchSupport.dispatchWithFailureCleanup(
                task, () -> invocationService.dispatch(task), state::releaseSlot, log);
    }
    if (state.queued() > 0 && dispatched > 0) {
        signalWork(functionName);
    }
}
\end{lstlisting}

Ogni visita a una funzione dispaccia al massimo due compiti, poi la funzione
viene rimessa in fondo alla coda dei nomi attivi se ha ancora lavoro. \`E un
round-robin con quanto pari a due.

Il quanto limitato \`e ci\`o che garantisce l'equit\`a fra funzioni. Senza di
esso, una funzione con mille compiti accodati e slot disponibili terrebbe il
thread dello scheduler finch\'e non li avesse dispacciati tutti, e una funzione
con un solo compito accodata subito dopo attenderebbe l'intero drenaggio. Con il
quanto, la seconda attende al pi\`u due dispatch.

Il valore due, e non uno, riduce il costo di rotazione --- ogni ciclo comporta
un'operazione su due strutture concorrenti --- senza degradare
apprezzabilmente l'equit\`a. \`E un compromesso, e il progetto non ne documenta
una taratura sperimentale.

\begin{damisurare}
L'effetto di \code{MAX\_BATCH\_PER\_FUNCTION} sull'equit\`a fra funzioni
co-residenti e sul throughput aggregato non \`e stato misurato. Il
\cref{cap:25-valutazione-concorrenza} riporta misure di co-tenancy che
\emph{includono} l'effetto di questo parametro senza isolarlo.
\todomis{sensibilit\`a del throughput e della latenza al quanto di scheduling}
\end{damisurare}

\subsection{Pulizia in caso di fallimento del dispatch}

Il dispatch \`e affidato a un helper del nucleo, \code{dispatchWithFailureCleanup},
a cui viene passato il rilascio dello slot come azione di compensazione. Il
motivo \`e che lo slot \`e stato acquisito \emph{prima} del dispatch: se il
dispatch solleva prima di aver stabilito una catena di completamento, nessuno
rilascer\`a lo slot, e la funzione perderebbe permanentemente una unit\`a di
concorrenza. Ripetuto, il fenomeno porterebbe il limite effettivo a zero senza
che alcuna metrica lo segnali come anomalia --- il contatore delle richieste in
volo apparirebbe semplicemente sempre saturo.

\section{Le tre metriche esportate}

\code{QueueManager} registra tre gauge per ciascuna funzione al momento della
sua creazione:

\begin{center}
\small
\begin{adjustbox}{max width=\textwidth}
\begin{tabular}{@{}lll@{}}
\toprule
\textbf{Metrica} & \textbf{Sorgente} & \textbf{Consumatore} \\
\midrule
\code{function\_queue\_depth} & \code{state::queued} & autoscaler,
concurrency-control, operatore \\
\code{function\_inFlight} & \code{state::inFlight} & autoscaler,
concurrency-control \\
\code{function\_effective\_concurrency} & \code{state::effectiveConcurrency} &
osservazione della politica \\
\bottomrule
\end{tabular}
\end{adjustbox}
\end{center}

La terza \`e diagnostica e non alimenta alcuna decisione: rende osservabile
\emph{che cosa il controllore ha deciso}, distinguendolo da ci\`o che l'utente
aveva configurato. Senza di essa, osservare l'effetto di una politica adattiva
richiederebbe di dedurlo dal comportamento.

I meter sono rimossi quando la funzione viene rimossa, con la conseguenza gi\`a
discussa nel \cref{cap:06-nucleo} sugli scenari di verifica che interrogano le
metriche dopo la cancellazione.

\section{I due punti di estensione implementati}

\subsection{\code{QueueBackedEnqueuer}}

Implementa \code{InvocationEnqueuer} delegando a \code{QueueManager}. La sua
esistenza \`e ci\`o che fa passare \code{enabled()} da \code{false} a
\code{true}, e quindi ci\`o che fa rispondere \code{202} anzich\'e \code{501} a
\code{POST :enqueue}.

\subsection{\code{QueueBackedMetricsSource}}

Implementa \code{ScalingMetricsSource}, ed \`e la classe da cui dipende l'intera
Parte III di questo documento. Fornisce quattro operazioni: leggere la
profondit\`a della coda, leggere le richieste in volo, \emph{scrivere} il limite
di concorrenza effettivo, e registrare il modo del controllore attivo.

La terza \`e quella che conta. Un modulo di controllo della concorrenza calcola
un limite; \`e la scrittura attraverso questa interfaccia a rendere quel limite
effettivo, perch\'e da quel momento \code{tryAcquireSlot} lo confronta con il
contatore delle richieste in volo. \`E letteralmente l'unico punto in cui una
politica di concorrenza tocca il comportamento del sistema, e questo spiega
perch\'e \code{concurrency-control} rifiuti di avviarsi senza
\code{async-queue}: la scrittura verso un'implementazione no-op non fallisce, non
segnala nulla, e non ha effetto.

Il metodo \code{enabled()} distingue questa implementazione dal no-op del nucleo,
e nel modulo restituisce costantemente \code{true}.

\section{Il ruolo del modulo nel percorso sincrono}

Un'ultima nota, perch\'e il nome del modulo pu\`o trarre in inganno. La catena di
ammissione descritta nel \cref{cap:04-architettura-control-plane} usa la coda
asincrona anche per le invocazioni \emph{sincrone}, quando la coda sincrona non
\`e presente:

\begin{lstlisting}[language=Java]
private void admitLocally(ExecutionRecord executionRecord) {
    if (syncQueueGateway.enabled()) {
        syncQueueGateway.enqueueOrThrow(executionRecord.task());
    } else if (enqueuer.enabled()) {
        InvocationEnqueueSupport.enqueueOrThrow(enqueuer, metrics, executionRecord);
    } else {
        completionHandler.dispatch(executionRecord.task());
    }
}
\end{lstlisting}

In quella configurazione una richiesta sincrona viene accodata sulla coda della
funzione, il chiamante attende il futuro di completamento, e lo scheduler la
dispaccia quando uno slot si libera. Il risultato \`e che il limite di
concorrenza per funzione \`e applicato anche al traffico sincrono, il che \`e
esattamente ci\`o che si desidera. La differenza rispetto alla coda sincrona
(\cref{cap:10-sync-queue}) \`e che qui non esiste ammissione: una richiesta
sincrona che trova la coda piena riceve \code{429} da \code{QueueFullException}
senza che nessuno abbia stimato quanto avrebbe dovuto attendere.
